\documentclass[11pt,a4paper]{article}
\usepackage[utf8]{inputenc}
\usepackage[T1]{fontenc}
\usepackage[french]{babel}
\usepackage[margin=2.2cm]{geometry}
\usepackage{booktabs,longtable,array,tabularx}
\usepackage[table]{xcolor}
\usepackage{enumitem}
\usepackage{xurl}
\usepackage[hidelinks]{hyperref}
\usepackage{fancyhdr}

\definecolor{crit}{HTML}{B3261E}
\definecolor{maj}{HTML}{C76A00}
\definecolor{min}{HTML}{3A6EA5}
\definecolor{ok}{HTML}{2E7D32}
\newcommand{\crit}{\textcolor{crit}{\textbf{Critique}}}
\newcommand{\maj}{\textcolor{maj}{\textbf{Majeur}}}
\newcommand{\mineur}{\textcolor{min}{\textbf{Mineur}}}
\newcommand{\f}[1]{\nolinkurl{#1}}
\sloppy

\newenvironment{constat}[3]{%
  \subsection*{#1 --- #2}\noindent\textit{Gravité :} #3\par\smallskip
  \begin{description}[leftmargin=2.6cm,style=nextline,itemsep=2pt]}{\end{description}}

\pagestyle{fancy}
\fancyhf{}
\lhead{MyAnimLab --- Analyse assurance qualité}
\rhead{18 octobre 2026}
\cfoot{\thepage}

\title{\textbf{MyAnimLab}\\[4pt]\large Analyse d'assurance qualité du logiciel}
\author{Revue réalisée sur la branche \f{ccr-e82b7497-x6tqzs} (commit \f{a2b130b})}
\date{18 octobre 2026}

\begin{document}
\maketitle
\thispagestyle{fancy}

\section{Objet, périmètre et méthode}

MyAnimLab est une application web PHP/MySQL (dérivée de MyVivarium) de gestion
d'animalerie : registre d'entrées et de sorties, cages, lignées, essais,
protocoles et expérimentations, grilles cliniques et points limites,
autorisations éthiques, pharmacie et rapport annuel d'utilisation des animaux.
Ces données relèvent de la directive 2010/63/UE et de sa transposition française
(articles R.214-87 et suivants du Code rural, arrêtés du 1\textsuperscript{er}
février 2013). Leur exactitude et leur traçabilité sont donc une exigence
réglementaire, au-delà de la simple qualité logicielle.

\paragraph{Périmètre.} 101 scripts PHP à la racine, \f{lang/}, \f{tools/},
\f{setup/}, \f{docker/}, le schéma \f{database/schema.sql} et 26 migrations
(environ 26\,400 lignes PHP et 1\,900 lignes SQL).

\paragraph{Méthode.} Revue statique du code et de la configuration, contrôle
syntaxique (\texttt{php -l}) de tous les fichiers, recherches ciblées (requêtes SQL
construites par concaténation, contrôles d'authentification et de jeton CSRF,
téléversements, journalisation), vérification de la cohérence entre le schéma et
les migrations, mesure de la couverture du catalogue de traduction. L'application
n'a pas été exécutée : il n'y a pas eu de test dynamique ni de test d'intrusion.

\section{Synthèse}

\begin{center}
\rowcolors{2}{gray!8}{white}
\begin{tabularx}{\textwidth}{@{}lXc@{}}
\toprule
\textbf{Domaine} & \textbf{Constat principal} & \textbf{Évaluation}\\
\midrule
Sécurité & Fichier \f{.env} et téléversements servis par Apache ; quatre points d'entrée sans authentification & \textcolor{crit}{Insuffisant}\\
Intégrité des données & Modifier un suivi expérimental écrase les valeurs précédentes sans en garder trace & \textcolor{maj}{À renforcer}\\
Vérification & Aucun test automatisé, aucune intégration continue & \textcolor{crit}{Insuffisant}\\
Gestion de configuration & Migrations sans outil d'application ; la procédure d'installation déploie le dépôt amont & \textcolor{maj}{À renforcer}\\
Maintenabilité & Fichiers monolithiques, doublons \f{hc_*}/\f{bc_*}, contrôle CSRF recopié dans chaque page & \textcolor{maj}{À renforcer}\\
Accès aux données SQL & 357 requêtes préparées ; les rares interpolations portent sur des entiers & \textcolor{ok}{Satisfaisant}\\
Authentification & bcrypt, verrouillage après 3 échecs, régénération de l'identifiant de session & \textcolor{ok}{Satisfaisant}\\
Internationalisation & 1\,495 entrées ; tous les libellés \f{tr()} sont traduits & \textcolor{ok}{Satisfaisant}\\
\bottomrule
\end{tabularx}
\end{center}

\noindent Bilan : 3 constats critiques, 6 majeurs et 5 mineurs. Les constats
critiques sont des failles de sécurité exploitables dès le déploiement par
\f{setup/setup.sh}. Ils doivent être corrigés avant toute mise en production.

\section{Constats détaillés}

\subsection{Constats critiques}

\begin{constat}{QA-01}{Secrets et fichiers internes servis par le serveur web}{\crit}
\item[Preuve] \f{setup/setup.sh:149} définit \f{DocumentRoot} sur la racine du
  dépôt et le README demande d'y copier \f{.env}. Le dépôt ne contient aucun
  \f{.htaccess}, et la configuration Apache par défaut ne protège que les
  fichiers \f{.ht*}.
\item[Risque] Une simple requête \texttt{GET /.env} renvoie les identifiants MySQL et
  SMTP. \f{database/schema.sql}, \f{composer.lock}, \f{tools/} et \f{docs/} sont
  eux aussi publics.
\item[Action] Placer la racine web dans un sous-dossier \f{public/}, ou à défaut
  ajouter un \f{.htaccess} racine qui refuse \f{.env}, \f{*.sql}, \f{composer.*},
  \f{tools/}, \f{database/}, \f{setup/} et \f{vendor/}. Vérifier ensuite chaque
  déploiement avec \texttt{curl -I https://}\textit{<hôte>}\texttt{/.env}, qui doit répondre 403 ou 404.
\end{constat}

\begin{constat}{QA-02}{Points d'entrée sans authentification}{\crit}
\item[Preuve] \f{nt_edit.php} n'ouvre pas de session et modifie n'importe quelle
  note à partir de son identifiant. \f{nt_rmv.php} supprime n'importe quelle note
  sans vérifier la session ni l'auteur, alors que son en-tête affirme le
  contraire. \f{get_task.php} et \f{get_reminder.php} renvoient des données en
  JSON sans vérifier la session.
\item[Risque] Un visiteur anonyme peut altérer ou détruire des notes de cage
  (perte de traçabilité) et énumérer les tâches et rappels.
\item[Action] Créer un garde commun (\f{require_login()}, rôle et propriétaire),
  l'appliquer à ces quatre fichiers et ajouter le contrôle CSRF aux deux points
  d'entrée qui écrivent.
\end{constat}

\begin{constat}{QA-03}{Téléversements de cage publics et SVG accepté}{\crit}
\item[Preuve] \f{hc_edit.php}, l.~341--372 et \f{bc_edit.php}, l.~258--291 écrivent dans
  \f{uploads/}\textit{<cage>}\f{/}\textit{<nom d'origine>}, créé avec les droits \f{0777}. Les fichiers
  sont servis directement (\f{hc_view.php:507}). L'extension \f{svg} est
  autorisée et le type MIME réel n'est pas contrôlé.
\item[Risque] Un SVG contenant du JavaScript s'exécute dans l'origine de
  l'application (XSS stockée, vol de session). Comme les chemins sont
  prévisibles, les documents sont accessibles sans être connecté.
\item[Action] Reprendre le modèle de \f{ethics_lib.php} : stockage hors de la
  racine web ou derrière un \texttt{Require all denied}, nom aléatoire, contrôle par
  \f{finfo}, téléchargement par un script qui vérifie les droits. Retirer
  \f{svg} de la liste et passer les droits à \f{0750}.
\end{constat}

\subsection{Constats majeurs}

\begin{constat}{QA-04}{Écrasement silencieux des suivis expérimentaux}{\maj}
\item[Preuve] \f{experiment_followup.php}, l.~215--242 : une modification fait un
  \texttt{ON DUPLICATE KEY UPDATE} sur \f{experiment_followups}, puis un
  \f{DELETE} et une réinsertion des scores et des observations. L'ancienne valeur
  n'est conservée nulle part et la modification n'est pas journalisée.
\item[Risque] Les principes ALCOA+ (données attribuables, originales, durables)
  ne sont plus respectés. Lors d'une inspection DDPP ou d'une revue de comité
  d'éthique, on ne peut pas reconstituer le score clinique qui a motivé, ou non,
  l'application d'un point limite.
\item[Action] Versionner les suivis (table d'historique, ou ligne \og
  remplacée\fg{} avec un \f{superseded_by}), exiger un motif de modification et
  journaliser l'avant et l'après dans \f{activity_log}. Appliquer le même
  principe aux autres écrans de modification des données réglementaires
  (registre, sorties, prescriptions).
\end{constat}

\begin{constat}{QA-05}{Journal d'audit modifiable et incomplet}{\maj}
\item[Preuve] La table \f{activity_log} (\f{schema.sql:792}) n'enregistre que du
  texte libre, sans valeurs avant ni après. Le compte MySQL de l'application peut
  y faire \f{UPDATE} et \f{DELETE}.
\item[Action] Créer un compte MySQL limité à \f{INSERT} et \f{SELECT} sur les
  tables d'audit, ajouter des champs structurés (\f{old_value},
  \f{new_value}, motif) et, si besoin, un chaînage par empreinte des
  enregistrements pour rendre toute altération détectable.
\end{constat}

\begin{constat}{QA-06}{Aucune vérification automatisée}{\maj}
\item[Preuve] Le dépôt ne contient ni répertoire de tests, ni PHPUnit dans
  \f{composer.json}, ni workflow dans \f{.github/}, ni outil d'analyse statique.
  Les 112 commits (dont 64 générés par un assistant) ont été intégrés sans
  contrôle automatique, et seules 2 pull requests ont été fusionnées.
\item[Action] Mettre en place un socle minimal :
  \begin{enumerate}[nosep]
    \item une CI GitHub Actions avec \texttt{php -l}, PHPStan (niveau 5 au départ),
      \texttt{composer audit} et le chargement de \f{schema.sql} dans MySQL~8 ;
    \item des tests PHPUnit des fonctions métier de \f{animal_lib.php} et
      \f{ethics_lib.php} (validité des autorisations, sévérité, calcul du
      rapport annuel) ;
    \item quelques tests de bout en bout Playwright sur les parcours réglementaires
      (entrée d'un animal, suivi, point limite, sortie, rapport annuel) ;
    \item une règle de protection de branche avec revue humaine obligatoire.
  \end{enumerate}
\end{constat}

\begin{constat}{QA-07}{Gestion du schéma et du déploiement}{\maj}
\item[Preuve] Les 26 migrations de \f{database/migrations/} n'ont ni outil
  d'application ni table de version, et \f{schema.sql} est tenu à jour à la main
  en parallèle (il est cohérent à ce jour). Par ailleurs, \f{setup/setup.sh:4}
  et le README clonent \f{myvivarium/MyVivarium}, c'est-à-dire le dépôt amont et
  non ce fork.
\item[Risque] Une installation \og en un clic\fg{} déploie une autre version du
  logiciel, et les mises à jour en production risquent d'oublier ou de rejouer
  des migrations.
\item[Action] Ajouter une table \f{schema_migrations} et un script
  \f{tools/migrate.php} idempotent. En CI, vérifier que l'application des
  migrations à partir de \f{001} produit le même schéma que \f{schema.sql}.
  Corriger les URL du README et de \f{setup.sh}.
\end{constat}

\begin{constat}{QA-08}{Scripts cron accessibles par HTTP}{\maj}
\item[Preuve] \f{process_reminders.php} ne vérifie pas \f{PHP_SAPI}, alors que
  \f{daily_round_alert.php:14} le fait. Il suffit donc d'appeler son URL pour
  générer des tâches et remplir la file d'e-mails.
\item[Action] Ajouter le même garde CLI et, plus largement, sortir les scripts
  cron de la racine web.
\end{constat}

\begin{constat}{QA-09}{Export complet contenant des secrets}{\maj}
\item[Preuve] \f{export_for_v2.php} exporte en JSON la table \f{users} entière,
  y compris les hachés de mots de passe et les jetons de réinitialisation.
  L'export est réservé aux administrateurs.
\item[Action] Exclure les colonnes sensibles, journaliser chaque export et
  limiter la durée de vie de la fonction à la période de migration.
\end{constat}

\subsection{Constats mineurs}

\begin{constat}{QA-10}{Maintenabilité}{\mineur}
\item[Preuve] Plusieurs fichiers dépassent 1\,000 lignes (\f{hc_edit.php},
  \f{bc_edit.php}) et les couples \f{hc_*} / \f{bc_*} sont presque identiques.
  Logique métier, SQL et HTML sont mêlés dans les mêmes fichiers. Le contrôle
  CSRF est recopié dans 44 fichiers sans fonction commune.
\item[Action] Factoriser progressivement : fonctions \f{csrf_field()} et
  \f{csrf_verify()}, repositories par entité, un seul écran paramétré pour les
  deux catégories de cages.
\end{constat}

\begin{constat}{QA-11}{Fuite d'informations techniques}{\mineur}
\item[Preuve] \f{nt_add.php} et \f{nt_rmv.php} renvoient \f{$con->error} et
  \f{$stmt->error} au client.
\item[Action] Journaliser l'erreur côté serveur et renvoyer un message générique.
\end{constat}

\begin{constat}{QA-12}{Requêtes SQL interpolées}{\mineur}
\item[Preuve] Une dizaine de requêtes interpolent des listes d'identifiants
  (\f{experiment_report.php:63}, \f{litter_manage.php:170}, etc.). Toutes les
  valeurs sont converties en entiers, donc sans risque aujourd'hui.
\item[Action] Utiliser partout des requêtes préparées, pour qu'une évolution
  future ne puisse pas introduire d'injection.
\end{constat}

\begin{constat}{QA-13}{Valeurs par défaut inadaptées}{\mineur}
\item[Preuve] Le fuseau horaire par défaut de \f{process_reminders.php} est
  \f{America/New_York}. Si le paramètre \f{timezone} est absent, les rappels et
  les tâches sont décalés.
\item[Action] Utiliser \f{Europe/Paris} par défaut, ou refuser de s'exécuter
  tant que le paramètre n'est pas défini.
\end{constat}

\begin{constat}{QA-14}{Hygiène du dépôt}{\mineur}
\item[Preuve] \f{images/.DS_Store} est versionné et \f{uploads/} n'est pas exclu
  dans \f{.gitignore}. Les contraintes de version de PHP et des extensions
  (\f{mysqli}, \f{fileinfo}) ne sont pas déclarées dans \f{composer.json}.
\item[Action] Nettoyer le dépôt, compléter \f{.gitignore}, déclarer
  \texttt{"php": ">=8.1"} et \f{ext-*} dans \f{composer.json}.
\end{constat}

\section{Points forts}

\begin{itemize}[nosep]
  \item 357 requêtes préparées. Aucune injection SQL n'a été trouvée.
  \item Mots de passe hachés avec bcrypt, compte verrouillé 15 minutes après
    3 échecs, identifiant de session régénéré à la connexion, cookies
    \f{HttpOnly}, \f{Secure} et \f{SameSite=Strict}.
  \item Documents d'éthique bien protégés (\f{ethics_lib.php}) : contrôle MIME,
    nom aléatoire, accès direct refusé, taille limitée.
  \item Suppression définitive d'un animal réservée aux administrateurs,
    confirmée par saisie de l'identifiant et journalisée avant suppression.
  \item Journaux métier spécialisés (événements et mouvements d'animaux, cycle de
    vie des cages, portées), transactions dans 18 scripts, 74 clés étrangères.
  \item Contrôle syntaxique réussi sur 100\,\% des fichiers PHP, catalogue de
    traduction complet, environnement Docker local isolé de la production.
\end{itemize}

\section{Plan d'action priorisé}

\begin{center}
\rowcolors{2}{gray!8}{white}
\begin{tabularx}{\textwidth}{@{}clXl@{}}
\toprule
\textbf{Priorité} & \textbf{Constats} & \textbf{Action} & \textbf{Effort}\\
\midrule
P0 & QA-01, 02, 03, 08 & Racine web dans \f{public/}, garde d'authentification commun, téléversements protégés, garde CLI & 2--3 j\\
P1 & QA-04, 05 & Versionnement des suivis, motif obligatoire, audit avant et après non modifiable & 3--5 j\\
P1 & QA-06 & CI (lint, PHPStan, \texttt{composer audit}, schéma), premiers tests PHPUnit, protection de branche & 3 j\\
P2 & QA-07, 09 & Outil de migrations, URL du dépôt, export sans secrets & 2 j\\
P3 & QA-10 à 14 & Factorisation, nettoyage, valeurs par défaut & continu\\
\bottomrule
\end{tabularx}
\end{center}

\noindent Pour un usage réglementaire, nous recommandons aussi de rédiger une
courte documentation de validation : exigences, matrice de traçabilité entre
exigences et tests, et procédure de gestion des changements. Elle sert de
preuve de maîtrise du logiciel lors des inspections.

\appendix
\section{Métriques relevées}

\begin{center}
\begin{tabular}{@{}lr@{}}
\toprule
Scripts PHP (racine) & 101\\
Lignes PHP (y compris \f{lang/} et \f{tools/}) & $\approx 26\,400$\\
Lignes SQL (schéma et migrations) & $\approx 1\,900$\\
Tables du schéma / clés étrangères & 55 / 74\\
Migrations & 26\\
Requêtes préparées & 357\\
Fichiers avec contrôle CSRF & 44\\
Entrées du catalogue français / libellés manquants & 1\,495 / 0\\
Tests automatisés / workflows CI & 0 / 0\\
Erreurs de syntaxe (\texttt{php -l}) & 0\\
\bottomrule
\end{tabular}
\end{center}

\end{document}
